%! Populations, Samples, Parameters, and Statistics — Unit 1, Lesson 1.2 — Guided Notes & Practice
% THE in-class document, in the MAIN IDEAS / NOTES shape (LESSON_SHAPE.md §1, ported from AP
% Statistics 2026-09-12): the vocabulary box, the hook, then ONE two-column guidednotes table.
% Each row is one idea — a label on the left; on the right one or two COMPLETE printed
% sentences, ONE large pre-drawn display the student reads or names, then a bold prompt and a
% two-across grid of numbered problems with 1.8–2.6cm of writing room. Problems run 1..19.
% DENSITY: a \blank{} only where a word IS the answer (the five cells of the constraint table
% in row 3 — the one \wordbank strip sits above it); no mid-sentence blanks; every display is
% pre-drawn; the page belongs to the student's pen.
%   I DO   : the teacher reads the definition, marks up the display, works the first problem
%            of each grid (1, 5, 10, 12).
%   WE DO  : the class works the rest of each grid, students holding the pen; the last row,
%            Guided Practice (16–19), is worked entirely together on Riverbend instead of
%            the market.
%   YOU DO : nothing on this page — ../homework/ IS the individual practice, started in the
%            last ten minutes of the period. No independent set, no reflection box, no
%            objectives box (the cover carries the targets).
% THE TRAP is problem 8 (61% from the records of all 900 — a percent that is a parameter).
% THE CRUX is problem 14 (row 4, "One Parameter"): three honest samples, three statistics,
% one parameter that never moved — the hook vote is re-taken there (PS.DC.1d–e). Problem 18
% is the crux transferred to Riverbend; problem 19 is the over-claim (an estimate is not a
% count).
%
% THE DATA
%   Lakeside Farmers Market: 800 Saturday shoppers; 50 asked, 26 drove -> 52%; 0.52 x 800 = 416
%   Three volunteers, a different 50 shoppers each: Ana 26 (52%), Ben 24 (48%), Cleo 29 (58%)
%   District records: 61% of all 900 Riverbend students ride a bus  (a PARAMETER — the trap)
%   Guided Practice, Riverbend HS: 900 students; 50 asked, 20 have a job -> 40%; 0.40 x 900 = 360;
%     a second 50 asked, 23 have one -> 46%
%
% PAGE LOCKSTEP with notes_key/: the key is GENERATED from this file by templates/lesson/mkkey.py
% (spec: templates/lesson/examples/lesson02_notes_key.json) — every \blank{} becomes an \ans{},
% every \vterm{} a \vtermans{}{}, every \par\writelines{n} n \ansline{} calls, and the answer
% argument of every \pcell, \writespace and \labelbox is filled. Answer spaces are fixed
% heights, so the two files cannot drift. KEEP EVERY \pcell ON ONE SOURCE LINE — mkkey matches
% line by line.
% TABLE MECHANICS: inside a cell, `\\` ENDS THE ROW — break lines with \par. A row cannot break
% across pages: the figure gets its own sub-row (`\\` then `& ...`) and EACH GRID ROW is its own
% sub-row — close the probgrid, `\\ &`, reopen as probgrid* (no top rule).
\documentclass[12pt]{article}
\usepackage{saar-article}
\usepackage{saar-boxes}
\newcommand{\TallMath}[1]{$\displaystyle #1\rule[-1.4em]{0pt}{3.2em}$}
% Word bank strip — ONLY above a table the student fills (row 3). Defined per document;
% the key inherits it and prints the same strip, so heights match.
\newcommand{\wordbank}[1]{\par\vspace{2pt}\noindent\fcolorbox{goldacc}{goldbg}{\parbox{\dimexpr\linewidth-2\fboxsep-2\fboxrule\relax}{\small\textbf{Word bank:} #1}}\par\vspace{4pt}}
% Vocabulary rows: a FIXED-HEIGHT row carrying the term label and OPEN writing space —
% no underline beside the term and no rule line (user correction 2026-09-06: the black
% answer line crowded the row; students write beside and below the term). The key fills
% exactly the same height, so the box cannot drift blank vs. keyed. saar-article's
% \termblank adds an inline blank and a rule line, so the pair is defined here (the key is
% generated from this file and inherits both). 2.3cm per row gives students three
% handwritten lines of room; keep key definitions to two lines.
\newcommand{\termrowheight}{2.3cm}
\newlength{\termrowinset}\setlength{\termrowinset}{7pt}
\newcommand{\vterm}[1]{%
  \par\noindent\begin{minipage}[t][\termrowheight][t]{\linewidth}%
    \vspace*{\termrowinset}%
    \textbf{\textcolor{cerulean}{#1:}}%
  \end{minipage}\par}
\newcommand{\vtermans}[2]{%
  \par\noindent\begin{minipage}[t][\termrowheight][t]{\linewidth}%
    \vspace*{\termrowinset}%
    \textbf{\textcolor{cerulean}{#1:}}\quad\ans{#2}%
  \end{minipage}\par}

\begin{document}

\pageheader{Unit 1, Lesson 1.2}{Guided Notes \& Practice}
% NAMESTRIP: no \namedateperiod here. The student writes their name once, on the
% cover this component is stapled behind.

\vspace{0.15in}

\begin{vocabbox}
\small
Fill in each term \textbf{as we name it} in the notes below.
\vterm{Population}
\vterm{Sample}
\vterm{Parameter}
\vterm{Statistic}
\vterm{Sampling variability}
\end{vocabbox}

\boxguard[12]
\begin{hookbox}
\small
Three volunteers stood at the Lakeside Farmers Market on the \emph{same} Saturday morning, and
each asked a \emph{different} $50$ shoppers whether they drove. Ana found $52\%$. Ben found
$48\%$. Cleo found $58\%$.

\vspace{3pt}
\textbf{Who made the mistake?} Circle one: \quad Ana \qquad Ben \qquad Cleo \qquad nobody
\quad --- and say why you think so. We vote again at problem~14.
\par\writelines{2}
\end{hookbox}

\small
\begin{guidednotes}
% ── ROW 1: the two groups ────────────────────────────────────────────────────
\mainidea[Two groups]{Population \& Sample} &
The \textbf{population} is the entire group the question is about --- everyone the conclusion
should describe. The \textbf{sample} is the part you actually collected data from; its
\textbf{sample size} is $n$. \textbf{The population is decided by the question, not by who was handy.}
\\
& {\centering
\begin{tikzpicture}[scale=0.72]
  % the population: a box of 800 shoppers (each dot stands for a few)
  \fill[frost, rounded corners=6pt] (-0.5,-0.55) rectangle (10.5,3.1);
  \draw[cerulean, thick, rounded corners=6pt] (-0.5,-0.55) rectangle (10.5,3.1);
  \node[anchor=west, font=\small\bfseries\color{cerulean}] at (-0.3,2.72)
    {Lakeside Farmers Market, Saturday morning --- $800$ shoppers};
  % the sample: an oval of 50, drawn under the dots
  \fill[goldbg] (3.5,1.0) ellipse (1.55 and 0.95);
  \draw[goldacc, very thick] (3.5,1.0) ellipse (1.55 and 0.95);
  \foreach \i in {0,...,20} { \foreach \j in {0,...,4} {
    \fill[slate] ({\i*0.5},{\j*0.5}) circle (0.075); }}
  \node[font=\scriptsize\bfseries\color{goldacc}, fill=goldbg, inner sep=1.5pt] at (3.5,2.15)
    {$50$ asked --- $26$ drove};
  \node[font=\tiny, anchor=east] at (10.35,-0.3) {\emph{each dot: a few shoppers}};
\end{tikzpicture}\par\vspace{5pt}
Box: the \labelbox{2.2cm}{}\quad Oval: the \labelbox{1.8cm}{}\quad $n = $ \labelbox{1cm}{}\par}
\\
& \notesprompt{Name the two groups and the sample size --- \emph{the population first; the sample is chosen out of it.}}
\begin{probgrid}{|Y|Y|}
\pcell{1}{The manager asks $50$ of the market's $800$ Saturday shoppers how they traveled.}{1.6cm}{} & \pcell{2}{The principal asks $45$ of Riverbend's $900$ students how many hours a week they work.}{1.6cm}{} \\ \hline
\end{probgrid}
\\
& \begin{probgrid}*{|Y|Y|}
\pcell{3}{A cereal plant weighs $30$ of the $12{,}000$ boxes it filled on Tuesday.}{1.6cm}{} & \pcell{4}{Data from \emph{every} individual is a \textbf{census}. Suppose the manager asked all $800$. Name the population and the sample now.}{2.0cm}{} \\ \hline
\end{probgrid}
\\ \hline
% ── ROW 2: the two kinds of numbers ──────────────────────────────────────────
\mainidea[Two kinds of numbers]{Parameter \& Statistic} &
A number that describes the \textbf{population} is a \textbf{parameter}. A number computed
from the \textbf{sample} is a \textbf{statistic}. \textit{Memory hook:}
\textbf{p}opulation--\textbf{p}arameter, \textbf{s}ample--\textbf{s}tatistic. To sort a
number, ask one question: \textbf{who does it describe?}
\\
& {\centering
\begin{tikzpicture}[scale=0.8]
  % left panel: the sample and its statistic
  \fill[goldbg, rounded corners=5pt] (0,0) rectangle (4.8,2.9);
  \draw[goldacc, thick, rounded corners=5pt] (0,0) rectangle (4.8,2.9);
  \node[font=\footnotesize\bfseries\color{goldacc}] at (2.4,2.5) {SAMPLE --- the $50$ asked};
  \node[font=\scriptsize] at (2.4,2.0) {$26$ of them drove};
  \node[font=\normalsize\bfseries] at (2.4,1.15) {$\dfrac{26}{50} = 52\%$};
  \node[font=\tiny\itshape] at (2.4,0.35) {known --- you counted it};
  % right panel: the population and its parameter
  \fill[frost, rounded corners=5pt] (6.8,0) rectangle (11.6,2.9);
  \draw[cerulean, thick, rounded corners=5pt] (6.8,0) rectangle (11.6,2.9);
  \node[font=\footnotesize\bfseries\color{cerulean}] at (9.2,2.5) {POPULATION --- all $800$};
  \node[font=\scriptsize] at (9.2,2.0) {the true percent who drove};
  \node[font=\normalsize\bfseries] at (9.2,1.15) {$?\,\%$};
  \node[font=\tiny\itshape] at (9.2,0.35) {unknown --- and there is only one};
  % the arrow between them
  \draw[->, very thick, charcoal] (4.9,1.3) -- (6.7,1.3);
  \node[font=\tiny\bfseries, above] at (5.8,1.35) {estimates};
\end{tikzpicture}\par\vspace{5pt}
The $52\%$ is a \labelbox{2.4cm}{}\qquad the true percent for all $800$ is a \labelbox{2.4cm}{}\par}
You \emph{have} the statistic; you \emph{want} the parameter; you use the one to \textbf{estimate} the other.
\\
& \notesprompt{Parameter or statistic? Say who the number describes.}
\begin{probgrid}{|Y|Y|}
\pcell{5}{$52\%$ of the $50$ shoppers asked drove to the market.}{1.2cm}{} & \pcell{6}{The market had $800$ shoppers on Saturday morning.}{1.6cm}{} \\ \hline
\end{probgrid}
\\
& \begin{probgrid}*{|Y|Y|}
\pcell{7}{The $8$ shoppers asked traveled a mean of $4$ miles.}{1.2cm}{} & \pcell{8}{District records show $61\%$ of Riverbend's $900$ students ride a bus. \textbf{Vote first, then decide.}}{1.6cm}{} \\ \hline
\end{probgrid}
\\ \hline
% ── ROW 3: constraints — why we sample at all ────────────────────────────────
\mainidea[Why we sample at all]{Constraints} &
If the manager wants a parameter, why not just ask all $800$? A \textbf{constraint} is a
real-world limit on how a study can be run --- the reason a census is usually off the table.
\\
& {\centering
\begin{tikzpicture}
  \foreach \i/\cname/\cdesc in {
      0/Time/{longer than\\you have},
      1/Cost/{more than\\the budget},
      2/Access/{can't be found,\\reached, or\\made to answer},
      3/Destruction/{measuring it\\ruins it},
      4/Change/{the population\\keeps changing}} {
    \fill[frost, rounded corners=4pt] ({\i*2.35},0) rectangle ({\i*2.35+2.2},1.7);
    \draw[cerulean, thick, rounded corners=4pt] ({\i*2.35},0) rectangle ({\i*2.35+2.2},1.7);
    \node[font=\small\bfseries\color{cerulean}] at ({\i*2.35+1.1},1.38) {\cname};
    \node[font=\tiny, align=center] at ({\i*2.35+1.1},0.62) {\cdesc};
  }
\end{tikzpicture}\par}
\\
& \textbf{9.} Name the constraint that stands in the way in each situation.
\footnotesize\renewcommand{\arraystretch}{1.15}
\begin{tabularx}{\linewidth}{@{} X p{2.4cm} @{}}
\toprule
\textbf{Situation} & \textbf{Constraint} \\
\midrule
Weighing a cereal box means opening it ($12{,}000$ filled).
  & \blank{1.9cm} \\
A different set of shoppers comes every Saturday.
  & \blank{1.9cm} \\
Asking all $900$ students takes three class periods.
  & \blank{1.9cm} \\
Reaching every county household means $40{,}000$ letters.
  & \blank{1.9cm} \\
Shoppers pay and leave before a clipboard reaches them.
  & \blank{1.9cm} \\
\bottomrule
\end{tabularx}\small
\\
& \notesprompt{A constraint is not an excuse --- it is why sampling exists.}
\begin{probgrid}{|Y|Y|}
\pcell{10}{One manager, one clipboard, one morning, $800$ shoppers. Which constraint kept her from a census?}{1.6cm}{} & \pcell{11}{What would she have to do to get the \emph{parameter} itself --- and why is that off the table?}{1.6cm}{} \\ \hline
\end{probgrid}
\\ \hline
% ── ROW 4: THE CRUX — three samples, one parameter ───────────────────────────
\mainidea[Three samples,]{One Parameter} &
Three volunteers each asked a \emph{different} $50$ shoppers on the \emph{same} Saturday
morning. Ana's $50$: $26$ drove. Ben's $50$: $24$ drove. Cleo's $50$: $29$ drove.
\textbf{The $800$ shoppers did not change between them.}
\\
& {\centering
\begin{tikzpicture}[scale=0.8]
  % the population with three samples inside it
  \fill[frost, rounded corners=6pt] (0,0) rectangle (5.2,3.2);
  \draw[cerulean, thick, rounded corners=6pt] (0,0) rectangle (5.2,3.2);
  \node[font=\scriptsize\bfseries\color{cerulean}] at (2.6,2.9) {all $800$ shoppers};
  \node[font=\scriptsize\color{cerulean}] at (2.6,2.55) {true percent who drove: $?$};
  \foreach \x/\y/\who/\c in {1.0/1.3/Ana/26, 2.6/0.8/Ben/24, 4.2/1.5/Cleo/29} {
    \fill[goldbg] (\x,\y) ellipse (0.8 and 0.62);
    \draw[goldacc, very thick] (\x,\y) ellipse (0.8 and 0.62);
    \node[font=\scriptsize\bfseries] at (\x,{\y+0.15}) {\who};
    \node[font=\tiny] at (\x,{\y-0.2}) {$\c$ of $50$};
  }
  % the number line of statistics
  \draw[->, thick] (6.2,1.2) -- (11.7,1.2);
  \foreach \p/\lab in {6.4/40, 7.6/45, 8.8/50, 10.0/55, 11.2/60} {
    \draw (\p,1.12) -- (\p,1.28); \node[below, font=\tiny] at (\p,1.08) {\lab\%}; }
  \foreach \p/\who in {8.32/Ben, 9.28/Ana, 10.72/Cleo} {
    \fill[goldacc] (\p,1.2) circle (0.11);
    \node[above, font=\tiny\bfseries] at (\p,1.35) {\who}; }
  \node[font=\scriptsize\bfseries, align=center] at (8.9,2.5) {three statistics};
  \node[font=\tiny\itshape\color{cerulean}, align=center] at (8.9,0.3)
    {the one parameter is somewhere here --- nobody knows where};
  \draw[->, goldacc, thick] (5.35,1.6) -- (6.25,1.4);
\end{tikzpicture}\par}
\\
& The change in a statistic from one sample to the next is called \textbf{sampling
variability}. It is not a mistake --- nobody miscounted; it is what samples do. \textbf{The
statistic moves; the parameter stays put.} Each statistic is an \textbf{estimate} of the one
parameter.\par
\textbf{In your own words:} why did three honest volunteers get three different numbers?
\writespace{1.4cm}{}
\\
& \notesprompt{Now settle the hook vote.}
\begin{probgrid}{|Y|Y|}
\pcell{12}{Compute Ben's and Cleo's statistics the way Ana's was computed: $26/50 = 0.52 = 52\%$.}{1.4cm}{} & \pcell{13}{How many statistics did the three volunteers produce? How many true percents are there for all $800$ shoppers?}{2.0cm}{} \\ \hline
\end{probgrid}
\\
& \begin{probgrid}*{|Y|Y|}
\pcell{14}{\textbf{Vote again:} Ana \; Ben \; Cleo \; nobody. Did one of them count wrong? Which one found the true percent for all $800$?}{2.0cm}{} & \pcell{15}{A fourth volunteer asks $200$ shoppers, not $50$. Would her statistic land closer to the true percent? Exactly on it?}{1.8cm}{} \\ \hline
\end{probgrid}
\\ \hline
% ── ROW 5: GUIDED PRACTICE — we do, together, a fresh context ────────────────
\mainidea[Guided practice]{Riverbend Jobs} &
The Riverbend principal wants to know what fraction of the school's $900$ students hold a
part-time job. She asks $50$ students; $20$ of them have one. The same week, the assistant
principal asks a \emph{different} $50$ students the same question; $23$ have one.
\\
& {\centering
\begin{tikzpicture}[scale=0.8]
  \fill[frost, rounded corners=6pt] (0,0) rectangle (11.2,2.7);
  \draw[cerulean, thick, rounded corners=6pt] (0,0) rectangle (11.2,2.7);
  \node[font=\small\bfseries\color{cerulean}] at (5.6,2.4) {Riverbend High School --- all $900$ students};
  \node[font=\scriptsize\color{cerulean}] at (5.6,2.05) {true percent with a part-time job: $?$};
  \foreach \x/\who/\c in {3.0/principal's $50$/20, 8.2/assistant principal's $50$/23} {
    \fill[goldbg] (\x,0.85) ellipse (2.1 and 0.65);
    \draw[goldacc, very thick] (\x,0.85) ellipse (2.1 and 0.65);
    \node[font=\scriptsize\bfseries] at (\x,1.05) {\who};
    \node[font=\scriptsize] at (\x,0.65) {$\c$ have a job};
  }
  \node[font=\large\bfseries\color{slate}] at (5.6,0.85) {vs.};
\end{tikzpicture}\par}
\\
& \notesprompt{We work these together.}
\begin{probgrid}{|Y|Y|}
\pcell{16}{Name the population, the sample, and $n$ --- and one constraint that kept her from asking all $900$.}{2.0cm}{} & \pcell{17}{What percent of those she asked have a job, and about how many of the $900$ is that? Label each number: statistic, parameter, or estimate.}{2.0cm}{} \\ \hline
\end{probgrid}
\\
& \begin{probgrid}*{|Y|Y|}
\pcell{18}{One survey says $40\%$; the other $46\%$. Did one make a mistake? Which survey found the true percent for all $900$? Explain.}{2.0cm}{} & \pcell{19}{Her draft report: \emph{``Exactly $360$ Riverbend students have a part-time job.''} Allowed? Rewrite it to claim only what her sample allows.}{2.0cm}{} \\ \hline
\end{probgrid}
\\ \hline
\end{guidednotes}

% ── THE NOTES END AT GUIDED PRACTICE ─────────────────────────────────────────
% There is deliberately no "you do" here: ../homework/ is the individual practice,
% started in class in the last ten minutes while the teacher circulates.

\end{document}
